\section{Symmetry of the Weyl frame: the Weil representation and its splitting problem}
\label{sec:weil}

\paragraph{Recollection.} $\kappa=\mathbb{F}_q$, $q=p^m$; $V(\kappa)=\kappa\oplus\kappa$
with $\omega\big((a,b),(a',b')\big)=ab'-a'b$ (Definition~\ref{def:symplectic-object});
$SL_2(\kappa)$ is the group of $\kappa$-linear automorphisms of $V(\kappa)$
preserving $\omega$ (Theorem~\ref{thm:wh-form}(e)); a polarizing cocycle
$\beta\in\mathrm{Adm}(\omega)$ and phase datum $\psi$ determine the algebra
$A_{\psi,\beta}(V)=\bigoplus_v\mathbb{C}\cdot W_\beta(v)$
(Definitions~\ref{def:admissible-cocycles-restated}, \ref{def:phase-datum},
\ref{def:weyl-algebra}), and $Q_\beta(v):=\beta(v,v)$
(Definition~\ref{def:quadratic-form}).
\provenance{D1, D2, D3, D4, D6}{---}{---}{---}

This section studies the symmetries of the Weyl frame itself: which
permutations of the Weyl operators extend to algebra automorphisms, how
large that symmetry group is, and whether $SL_2(\kappa)$ — the natural
symmetry group of $(V(\kappa),\omega)$ — genuinely acts on the quantum
system by honest algebra automorphisms (as opposed to merely projectively).
The answer is uniform in one respect and characteristic-dependent in
another: the group of frame automorphisms always surjects onto
$SL_2(\kappa)$ with kernel the dual group $\hat V$, at every characteristic;
but whether this extension \emph{splits} — whether $SL_2(\kappa)$ itself,
not just a projective shadow of it, acts by algebra automorphisms — is
settled affirmatively at odd $p$ and only partially resolved at $p=2$, where
it interacts with the Arf dichotomy of \S\ref{sec:polarizing-cocycle}.

\subsection{The Weyl frame and its two automorphism groups}

\begin{definition}[The Weyl frame, and its $\kappa$-linear automorphisms]
\label{def:weyl-frame}
The \emph{Weyl frame} is
\[
\mathcal{F} := \big\{\, \mathbb{C}^\times\cdot W_\beta(v) \ :\ v\in V(\kappa) \,\big\}
\ \subset\ A_{\psi,\beta}(V),
\]
the set of $\mathbb{C}^\times$-lines through Weyl operators. $\mathrm{Aut}_{\mathcal{F}}(A)$
is the group of $\mathbb{C}$-algebra automorphisms $\alpha$ of
$A_{\psi,\beta}(V)$ with $\alpha(\mathcal{F})=\mathcal{F}$.
$\mathrm{Aut}_{\mathcal{F}}^\kappa(A)\subseteq\mathrm{Aut}_{\mathcal{F}}(A)$ is
the subgroup consisting of those $\alpha$ whose induced permutation of
$V(\kappa)$ (well defined since $\alpha$ permutes the lines of $\mathcal{F}$,
which are indexed bijectively by $V(\kappa)$) is $\kappa$-linear.
\end{definition}

\begin{scope}
The $\kappa$-linearity clause in the definition of
$\mathrm{Aut}_{\mathcal{F}}^\kappa(A)$ is \textbf{imposed data}: nothing in
$A_{\psi,\beta}(V)$ or $\mathcal{F}$ as built from $(V,+)$ and the cocycle
$\psi\circ\beta$ forces an automorphism's induced permutation of $V(\kappa)$
to respect scalar multiplication by $\kappa$ rather than merely by
$\mathbb{F}_p$. Precisely how much symmetry this clause discards is
Theorem~\ref{thm:wh-symm} of \S\ref{sec:fp-symmetry}; nothing here
addresses that question.
\end{scope}

\provenance{D7}{---}{---}{---}

\subsection{Exactness of the frame-automorphism extension, uniformly in the
characteristic}

Fix $\beta\in\mathrm{Adm}(\omega)$. For $g\in SL_2(\kappa)$, write
$s_g(v,v') := \beta(gv,gv')-\beta(v,v')$ and $Q_g(v):=s_g(v,v)$; note that
$Q_g(v)=Q_\beta(gv)-Q_\beta(v)$.

\statusProved
\begin{theorem}[Exactness of the frame-automorphism extension]
\label{thm:wh-weil-a}
At every characteristic, the sequence
\[
1 \longrightarrow \hat V \longrightarrow \mathrm{Aut}_{\mathcal{F}}^\kappa(A_{\psi,\beta}(V))
\xrightarrow{\ \alpha\mapsto g\ } SL_2(\kappa) \longrightarrow 1
\]
is exact, where $\hat V := \mathrm{Hom}\big((V(\kappa),+),\mathbb{C}^\times\big)$; the
kernel consists exactly of the conjugations $\mathrm{Ad}_{W_\beta(u)}$,
$u\in V(\kappa)$.
\end{theorem}

\begin{scope}
This theorem does not assert that the extension splits — that $SL_2(\kappa)$
itself embeds back into $\mathrm{Aut}_{\mathcal{F}}^\kappa(A)$ as a
subgroup — only that it is exact. Splitting is addressed separately, and
with a genuinely different answer at odd and even characteristic, in
Theorems~\ref{thm:wh-weil-b} and \ref{conj:wh-weil-e} below. Nothing here
concerns $\mathrm{Aut}_{\mathcal{F}}(A)$ without the imposed $\kappa$-linearity
clause; that is Theorem~\ref{thm:wh-symm}.
\end{scope}

\provenance{D1, D4, D7, WH-ALG, WH-COMM, WH-WEIL-a}{theory/wh-kappa-choice.md \S8 \textit{step} $\langle 1\rangle 2,\langle 1\rangle 3$}{theory/checks/split24.py (X2)}{theory/verdicts/wh-kappa-adjudication.md}

\begin{proof}
\emph{$s_g$ is symmetric bi-additive, and $s_{gh}(v,v')=s_h(v,v')+s_g(hv,hv')$.}
Symmetry: $s_g(v,v')-s_g(v',v)=\omega(gv,gv')-\omega(v,v')=0$, since
$g\in SL_2(\kappa)$ preserves $\omega$ (Theorem~\ref{thm:wh-form}(e)); the
telescoping identity is immediate by expanding
$\beta(ghv,ghv')-\beta(v,v')$ through the intermediate term
$\beta(hv,hv')$. (For $\beta=\beta_0$ and $g=\left(\begin{smallmatrix}r&s\\t&u\end{smallmatrix}\right)$,
direct expansion of $\beta_0(gv,gv')=rt\,aa'+ru\,ab'+st\,a'b+su\,bb'$ gives
the explicit formula $s_g(v,v')=rt\,aa'+st(ab'+a'b)+su\,bb'$, using
$ru-st=\det g=1$.)

\emph{Description of $\mathrm{Aut}_{\mathcal{F}}^\kappa(A)$.} Frame
preservation gives $\alpha(W_\beta(v))=\lambda(v)W_\beta(gv)$ for some
function $\lambda:V(\kappa)\to\mathbb{C}^\times$ and some permutation $g$ of
$V(\kappa)$, imposed to be $\kappa$-linear by definition of
$\mathrm{Aut}_{\mathcal{F}}^\kappa(A)$. Applying $\alpha$ to
$W_\beta(v)W_\beta(v')=\psi(\beta(v,v'))W_\beta(v+v')$ and comparing basis
elements forces $g$ additive (automatic here, being $\kappa$-linear) and
\[
\lambda(v)\lambda(v')\,\psi\big(s_g(v,v')\big) = \lambda(v+v') \tag{$*$}
\]
for all $v,v'$. Applying $\alpha$ to the commutation relation of
Theorem~\ref{thm:wh-comm}(c) and using $\omega(gv,gv')=\det(g)\,\omega(v,v')$
(Theorem~\ref{thm:wh-form}(e)'s proof) gives
$\psi\big((\det g-1)\,\omega(v,v')\big)=1$ for all $v,v'$; since $\omega$ is
onto $\kappa$ (Theorem~\ref{thm:wh-form}(a),(c)), Fact~\ref{fact:val} forces
$\det g=1$, i.e.\ $g\in SL_2(\kappa)$. So $\alpha\mapsto g$ is a well-defined
map $\mathrm{Aut}_{\mathcal{F}}^\kappa(A)\to SL_2(\kappa)$, visibly a
homomorphism (composition of automorphisms composes the induced
permutations).

\emph{The kernel.} If $g=\mathrm{id}$, $(*)$ reads
$\lambda(v)\lambda(v')=\lambda(v+v')$, i.e.\ $\lambda\in\hat V$, a group of
order $q^2$ by Fact~\ref{fact:dual} ($U=V(\kappa)$, $d=2m$). The map
$u\mapsto\psi(\omega(u,\cdot))$, $V(\kappa)\to\hat V$, is injective (if
$\psi(\omega(u,\cdot))\equiv 1$ then $\omega(u,\cdot)\equiv 0$ by
Fact~\ref{fact:val}, so $u=0$ by nondegeneracy, Theorem~\ref{thm:wh-form}(c)),
hence bijective ($|V(\kappa)|=q^2=|\hat V|$); so every kernel element has
$\lambda=\psi(\omega(u,\cdot))$ for a unique $u\in V(\kappa)$, and by
Theorem~\ref{thm:wh-comm}(d) this is exactly $\mathrm{Ad}_{W_\beta(u)}$.

\emph{Surjectivity, uniformly in $p$.} For fixed $g\in SL_2(\kappa)$,
$c(v,v'):=\psi(s_g(v,v'))$ is a symmetric $2$-cocycle on the exponent-$p$
group $(V(\kappa),+)$ (symmetry of $c$ from symmetry of $s_g$, just shown;
the cocycle identity from bi-additivity of $s_g$), so Fact~\ref{fact:triv}
supplies $\lambda$ solving $(*)$: every $g\in SL_2(\kappa)$ is hit. This
argument is uniform in the characteristic; nothing distinguishes $p=2$.
\end{proof}

\begin{remark}[Every frame automorphism is projectively unitary, at every
characteristic]
\label{rem:proj-unitary}
Fix a model $(M,\pi)$ as in Theorem~\ref{thm:wh-svn}. Taking absolute values
in $(*)$ makes $v\mapsto|\lambda(v)|$ a homomorphism from the finite group
$V(\kappa)$ to the torsion-free group $(\mathbb{R}_{>0},\times)$, hence
trivial: every $\lambda$ occurring in
$\mathrm{Aut}_{\mathcal{F}}^\kappa(A)$ is automatically unimodular. By
Theorem~\ref{thm:wh-svn}(c),(d), every $\alpha\in\mathrm{Aut}_{\mathcal{F}}^\kappa(A)$
is $\mathrm{Ad}_T$ for some $T$, unique up to a scalar; since $\alpha$
permutes the \emph{unitary} Weyl operators (unimodular $\lambda$ times
unitary $\pi(W_\beta(gv))$), $T$ may be chosen unitary and is then unique up
to $U(1)$ (Theorem~\ref{thm:wh-svn}(d)). This gives an injective
homomorphism $\mathrm{Aut}_{\mathcal{F}}^\kappa(A)\to PU(M)$: a projective
unitary representation of $\mathrm{Aut}_{\mathcal{F}}^\kappa(A)$, at every
characteristic. Composed with a section $SL_2(\kappa)\to\mathrm{Aut}_{\mathcal{F}}^\kappa(A)$
of Theorem~\ref{thm:wh-weil-a}'s surjection — should one exist as an honest
group homomorphism, i.e.\ should the extension split — this yields an
honest (non-projective) representation of $SL_2(\kappa)$ by algebra
automorphisms, and hence, via $\mathrm{Ad}$, a projective unitary
representation of $SL_2(\kappa)$ itself on $M$. What is \emph{not} true is
that this extends to all of $\mathrm{Aut}_{\mathbb{C}}(A)\cong PGL_q(\mathbb{C})$:
at $q=2$, $T=\mathrm{diag}(2,1)$ conjugates $\pi(W(1,0))=\left(\begin{smallmatrix}0&1\\1&0\end{smallmatrix}\right)$
to $\left(\begin{smallmatrix}0&2\\1/2&0\end{smallmatrix}\right)$, whose singular values $2,\,1/2$ differ, so it is not a
scalar times a unitary — the frame hypothesis in
Definition~\ref{def:weyl-frame} is exactly what is needed above.
\end{remark}

\subsection{Splitting at odd characteristic}

\statusProvedConditional
\begin{theorem}[The Weil representation of $SL_2(\kappa)$, at odd characteristic]
\label{thm:wh-weil-b}
\textbf{[$p$ odd.]} $\alpha_g\big(W_\beta(v)\big) := \psi\big(Q_g(v)/2\big)\,W_\beta(gv)$
defines a splitting $g\mapsto\alpha_g$ of the extension of
Theorem~\ref{thm:wh-weil-a}: $SL_2(\kappa)$ acts on $A_{\psi,\beta}(V)$ by
$\mathbb{C}$-algebra automorphisms, with phases in $\mu_p$. Composed with
Remark~\ref{rem:proj-unitary}, this exhibits a projective unitary
representation of $SL_2(\kappa)$ on the model space of
Theorem~\ref{thm:wh-svn}.
\end{theorem}

\begin{scope}
This is an odd-characteristic statement only; Theorem~\ref{conj:wh-weil-e}
addresses whether an analogous honest (non-projective) action of
$SL_2(\kappa)$ by algebra automorphisms exists at $p=2$, and the answer
there is open in general. Nothing here claims the splitting is unique, or
canonical among splittings.
\end{scope}

\provenance{D4, D6, WH-WEIL-a, WH-SVN, WH-WEIL-b}{theory/wh-kappa-choice.md \S8 \textit{step} $\langle 1\rangle 4,\langle 1\rangle 7$}{theory/checks/split24.py (X4)}{theory/verdicts/wh-kappa-adjudication.md}

\begin{proof}
Since $s_g$ is symmetric bilinear, $Q_g(v+v')=Q_g(v)+Q_g(v')+2s_g(v,v')$
(the same polarization identity as in
Definition~\ref{def:quadratic-form}'s use elsewhere), so
\[
\psi\big(Q_g(v+v')/2\big) = \psi\big(Q_g(v)/2\big)\,\psi\big(Q_g(v')/2\big)\,\psi\big(s_g(v,v')\big),
\]
which is exactly relation $(*)$ of Theorem~\ref{thm:wh-weil-a}'s proof for
$\lambda(v):=\psi(Q_g(v)/2)$: $\alpha_g$ is a well-defined element of
$\mathrm{Aut}_{\mathcal{F}}^\kappa(A)$ over $g$. That $g\mapsto\alpha_g$ is a
homomorphism follows from the telescoping identity
$s_{gh}(v,v')=s_h(v,v')+s_g(hv,hv')$ of Theorem~\ref{thm:wh-weil-a}'s proof,
evaluated at $v=v'$ to give $Q_{gh}(v)=Q_h(v)+Q_g(hv)$, whence
$\alpha_g(\alpha_h(W_\beta(v))) = \psi(Q_h(v)/2)\,\psi(Q_g(hv)/2)\,W_\beta(ghv)
= \psi(Q_{gh}(v)/2)\,W_\beta(ghv) = \alpha_{gh}(W_\beta(v))$. This is an
honest algebra automorphism, so $SL_2(\kappa)$ acts by algebra automorphisms
with $\mu_p$ phases (division by $2$ makes sense since $p$ is odd); the
projective unitary statement is Remark~\ref{rem:proj-unitary} applied to
this section.
\end{proof}

\subsection{Characteristic two: the orthogonal subgroup and forced order-four
phases}

\statusProvedConditional
\begin{theorem}[$\mu_2$-phases occur exactly over the orthogonal subgroup of
$Q_\beta$]
\label{thm:wh-weil-c}
\textbf{[$p=2$.]} Write $O(Q_\beta) := \{\, g\in SL_2(\kappa) : Q_\beta\circ g = Q_\beta \,\}$.
A frame automorphism $\alpha_{g,\lambda}$ (in the notation
$\alpha(W_\beta(v))=\lambda(v)W_\beta(gv)$ of Theorem~\ref{thm:wh-weil-a}'s
proof) has \emph{some} choice of $\lambda$ landing entirely in
$\psi(\kappa)=\mu_2$ if and only if $g\in O(Q_\beta)$ — a condition on the
chosen $\beta$, not on the characteristic alone. For the reference cocycle
$\beta_0$, $O(Q_{\beta_0})\cap SL_2(\kappa)$ consists exactly of the matrices
$\mathrm{diag}(r,r^{-1})$ and $\mathrm{antidiag}(s,s^{-1})$, $r,s\in\kappa^\times$;
it has order $2(q-1)$ (equal to $2,6,14$ at $q=2,4,8$), and is a
\textbf{proper} subgroup of $SL_2(\kappa)$ for every $q$.
\end{theorem}

\begin{scope}
This theorem is about which $g$ admit a $\mu_2$-valued lift, for a
\emph{fixed} $\beta$; it says nothing yet about what happens for
$g\notin O(Q_\beta)$ (Theorem~\ref{thm:wh-weil-d}), and the concrete
description of $O(Q_{\beta_0})$ is specific to the reference cocycle — for a
general $\beta$ of hyperbolic type the order is the same, $2(q-1)$
(Theorem~\ref{thm:wh-beta-type}(iii)), but the concrete matrices differ with
the normal-form basis chosen there.
\end{scope}

\provenance{D6, WH-WEIL-a, D7, WH-WEIL-c}{theory/wh-kappa-choice.md \S8 \textit{step} $\langle 1\rangle 5$}{theory/checks/beta\_census.py (B6); theory/checks/fp\_symmetry.py (S3)}{theory/verdicts/wh-kappa-adjudication.md}

\begin{proof}
Setting $v'=v$ in $(*)$ at $p=2$ ($2v=0$) gives $\lambda(v)^2=\psi(Q_g(v))$.
So a solution $\lambda$ of $(*)$ is $\mu_2$-valued (i.e.\ $\lambda(v)^2=1$
for every $v$) iff $Q_g(V(\kappa))\subseteq\ker\psi$, i.e.\ iff
$\psi\circ Q_g\equiv 1$.

\emph{For every $\beta$, this is already the exact isometry condition.}
$Q_g(v)=Q_\beta(gv)-Q_\beta(v)$ is additive in $v$: expanding
$Q_\beta(g(v+v'))-Q_\beta(v+v')$ via Definition~\ref{def:quadratic-form}'s
polarization identity on both terms, the $\omega$-cross-terms cancel because
$g$ preserves $\omega$ (Theorem~\ref{thm:wh-form}(e)), leaving
$Q_g(v+v')=Q_g(v)+Q_g(v')$; also $Q_g(cv)=c^2Q_g(v)$ since $g$ is
$\kappa$-linear. So, in a basis, $Q_g(a,b)=\alpha a^2+\delta b^2$ for some
$\alpha,\delta\in\kappa$ (no cross term). If $\psi\circ Q_g\equiv 1$ then
$\psi(\alpha a^2)=1$ for all $a$ (setting $b=0$); since $a\mapsto a^2$ is
bijective on $\kappa$, this forces $\alpha=0$ by Fact~\ref{fact:val}, and
likewise $\delta=0$: $Q_g\equiv 0$, i.e.\ $g\in O(Q_\beta)$. Conversely, if
$g\in O(Q_\beta)$ then trivially $\psi\circ Q_g\equiv 1$, and moreover a
$\mu_2$-valued $\lambda$ solving $(*)$ then \emph{exists}: $\psi\circ s_g$ is
a symmetric cocycle with trivial diagonal ($\psi(s_g(v,v))=\psi(Q_g(v))=1$),
exactly the case handled by the $\ell$-construction in the proof of
Theorem~\ref{thm:wh-beta-level} (\S\ref{sec:polarizing-cocycle}). So
$\psi\circ Q_g\equiv 1$ iff $g\in O(Q_\beta)$ iff a $\mu_2$-valued lift
exists.

\emph{The reference cocycle.} By the explicit formula of
Theorem~\ref{thm:wh-weil-a}'s proof, $s_g(v,v')=rt\,aa'+st(ab'+a'b)+su\,bb'$
for $\beta_0$, so $Q_g(a,b)=rt\,a^2+su\,b^2$ (the cross term drops at
$p=2$). Thus $g\in O(Q_{\beta_0})$ iff $rt=0$ and $su=0$. Combined with
$\det g=ru-st=1$: if $r=0$ then $st=-1=1$ (so $s,t\neq0$), and $su=0$ forces
$u=0$, giving $g=\left(\begin{smallmatrix}0&s\\s^{-1}&0\end{smallmatrix}\right)$
(using $t=s^{-1}$ from $st=1$) — an antidiagonal matrix; if $t=0$ then
$ru=1$ (so $r,u\neq0$, $u=r^{-1}$), and $su=0$ forces $s=0$, giving
$g=\mathrm{diag}(r,r^{-1})$; the case $r=t=0$ contradicts $\det g=1$. So
$O(Q_{\beta_0})\cap SL_2(\kappa)=\{\mathrm{diag}(r,r^{-1})\}\cup\{\mathrm{antidiag}(s,s^{-1})\}$,
of order $(q-1)+(q-1)=2(q-1)$ (each family is parametrized bijectively by
$\kappa^\times$). Since $Q_{\beta_0}(a,b)=ab$ is the hyperbolic normal form
(Theorem~\ref{thm:wh-beta-type}(iii)), this matches the general formula
$2(q-1)$ there. Properness: $2(q-1) < q(q^2-1) = |SL_2(\kappa)|$ for every
$q\geq 2$, since $q(q+1)\geq 6 > 2$.
\end{proof}

\statusProvedConditional
\begin{theorem}[$\mu_4$ is forced exactly when the orthogonal subgroup is
proper]
\label{thm:wh-weil-d}
\textbf{[$p=2$.]} If $g\notin O(Q_\beta)$, \emph{every} frame automorphism
over $g$ carries a phase of exact multiplicative order $4$. Consequently
$\mu_4$ is forced somewhere in $\mathrm{Aut}_{\mathcal{F}}^\kappa(A)$ exactly
when $O(Q_\beta)\cap SL_2(\kappa)$ is a proper subgroup of $SL_2(\kappa)$:
this holds for $\beta_0$ at every $q$ (Theorem~\ref{thm:wh-weil-c}), and for
\emph{every} $\beta\in\mathrm{Adm}(\omega)$ at $q=4,8$, by direct
enumeration. It does \textbf{not} hold in the single case $q=2$ with $\beta$
anisotropic, where $O(Q_\beta)\cap SL_2(\mathbb{F}_2)=SL_2(\mathbb{F}_2)$
(order $6=|O(Q_\beta)\cap SL_2(\kappa)|$ from the anisotropic formula
$2(q+1)$ of Theorem~\ref{thm:wh-beta-type}(iii) at $q=2$) and the entire
group $\mathrm{Aut}_{\mathcal{F}}^\kappa(A)$ admits $\mu_2$-valued phases
throughout. So $\mu_4$ is \textbf{not} a characteristic-two invariant of
$(V(\kappa),\omega,\psi)$ alone: whether it is forced depends on the type of
$\beta$.
\end{theorem}

\begin{scope}
This theorem uses only Theorem~\ref{thm:wh-weil-c} and the comparison of
$2(q\mp1)$ against $|SL_2(\kappa)|=q(q^2-1)$; it does not itself re-derive
the general-$q$ order $2(q+1)$ for the anisotropic stabilizer (that is
Theorem~\ref{thm:wh-beta-type}(iii), only cited here). The claim for
\emph{every} $\beta$ at $q=4,8$ rests on direct enumeration rather than a
closed-form argument for general $q$.
\end{scope}

\provenance{D6, D10, WH-WEIL-c, WH-WEIL-d}{theory/wh-kappa-choice.md \S8 \textit{step} $\langle 1\rangle 6$}{theory/checks/split24.py (X4,X6, and the anisotropic cross-run at $q=2$)}{theory/verdicts/wh-kappa-adjudication.md}

\begin{proof}
By Theorem~\ref{thm:wh-weil-c}'s proof, any $\lambda$ solving $(*)$ over $g$
satisfies $\lambda(v)^2=\psi(Q_g(v))$ for every $v$. If $g\notin O(Q_\beta)$
then $Q_g\not\equiv 0$, and (by the same bijective-squaring argument as
there) $\psi(Q_g(v))=-1$ for some $v$ — since $Q_g(v)=-1$'s preimage under
$\psi$ being nontrivial requires only that $\psi(\kappa)=\mu_2=\{1,-1\}$
(Fact~\ref{fact:val}) and $Q_g$ nonzero, hence onto a nonzero multiple of
$\kappa$, hits every value including one where $\psi=-1$. Then
$\lambda(v)^2=-1$, i.e.\ $\lambda(v)$ has exact order $4$: \emph{every}
$\lambda$ solving $(*)$ over this $g$ (not just one choice) is forced to
take an order-$4$ value at this particular $v$, since the equation
$\lambda(v)^2=\psi(Q_g(v))=-1$ has no solution in $\mu_2$. Comparing
$2(q-1)$ and $2(q+1)$ (Theorem~\ref{thm:wh-beta-type}(iii)) against
$|SL_2(\kappa)|=q(q^2-1)=q(q-1)(q+1)$: $2(q-1)<q(q-1)(q+1)$ always
(as $q(q+1)\geq 6$ for $q\geq 2$), so the hyperbolic stabilizer — in
particular $O(Q_{\beta_0})\cap SL_2(\kappa)$ — is always proper, giving some
$g\notin O(Q_\beta)$ and hence a forced order-$4$ phase, at every $q$. For
the anisotropic stabilizer, $2(q+1)<q(q-1)(q+1)$ iff $2<q(q-1)$: this holds
for $q\geq 4$ but \emph{fails} at $q=2$, where $q(q-1)=2$ exactly, so
$2(q+1)=6=q(q-1)(q+1)=|SL_2(\mathbb{F}_2)|$: the anisotropic stabiliser is
the whole group at $q=2$, so \emph{no} $g$ lies outside $O(Q_\beta)$, so no
order-$4$ phase is forced by this argument, and — since $\mu_2$-valued lifts
exist for every $g\in O(Q_\beta)$ by Theorem~\ref{thm:wh-weil-c}'s proof, and
$\hat V$ is automatically $\mu_2$-valued (Fact~\ref{fact:dual}, exponent $p=2$)
— the whole of $\mathrm{Aut}_{\mathcal{F}}^\kappa(A)$ admits $\mu_2$-valued
phases in this one case. At $q=4,8$, exhaustive enumeration over every
$\beta\in\mathrm{Adm}(\omega)$ (\texttt{theory\slash{}checks\slash{}split24.py}, gates X4
and X6) confirms $O(Q_\beta)\cap SL_2(\kappa)$ is proper for both types,
hence $\mu_4$ is forced for every $\beta$ there.
\end{proof}

\subsection{Does the extension split at characteristic two? A conjecture, and
what is known}

\statusConjecture
\begin{conjecture}[Existence of a Weil representation of $SL_2(\kappa)$ at
characteristic two]
\label{conj:wh-weil-e}
\textbf{[$p=2$.]} The extension of Theorem~\ref{thm:wh-weil-a} splits for
every $q=2^m$: $SL_2(\kappa)$ acts on $A_{\psi,\beta}(V)$ by
$\mathbb{C}$-algebra automorphisms, for some choice of splitting. This is
verified for the reference cocycle $\beta_0$ at $q=2$ (exhaustive search
finds exactly $4$ distinct complements of $\hat V$ of order $6$, over all
generating pairs) and at $q=4$ ($64$ lifts of a $(2,3,5)$-presentation of
$SL_2(\mathbb{F}_4)$ satisfy the defining relations and generate a subgroup
of order $60$ meeting $\hat V$ trivially); in both verified cases every
complement necessarily uses phases of exact order $4$ (no complement avoids
$\mu_4$, since $|\hat V|\cdot|O(Q_{\beta_0})\cap SL_2(\kappa)|$ does not
divide $|\mathrm{Aut}_{\mathcal{F}}^\kappa(A)|$'s $\mu_2$-part suitably: $6\nmid 8$
at $q=2$, $60\nmid 96$ at $q=4$). No argument is known for general $q$.
\end{conjecture}

\begin{scope}
This is recorded as a conjecture, not a theorem, precisely because the
evidence is two finite instances and no general argument: per this
campaign's discipline, a passing computational instance promotes nothing.
The verified cases are for the reference cocycle $\beta_0$ only; splitting
for other cocycles, and any relationship between splittings for cocycles of
different Arf type, is not addressed. Should a splitting be found for
general $q$, Remark~\ref{rem:proj-unitary} would then give a projective
unitary representation of $SL_2(\kappa)$ at $p=2$ exactly as
Theorem~\ref{thm:wh-weil-b} does at odd $p$.
\end{scope}

\provenance{D7, WH-WEIL-a, WH-WEIL-d, WH-WEIL-e}{---}{theory/checks/split24.py (X4,X5,X6)}{theory/verdicts/wh-kappa-adjudication.md}

\begin{quote}
\emph{What is known, stated as evidence and not as a proof.} At $q=2$,
$|\mathrm{Aut}_{\mathcal{F}}^\kappa(A)|=q^2|SL_2(\kappa)|=4\cdot 6=24$;
exhaustive search over generating pairs of order-$6$ subgroups meeting
$\hat V=\mathbb{Z}/2\times\mathbb{Z}/2$ trivially finds exactly $4$
complements, all using a phase of exact order $4$ somewhere
(\texttt{split24.py} gate X4). At $q=4$, lifting a $(2,3,5)$
generator-relation presentation of $SL_2(\mathbb{F}_4)\cong A_5$ over all
$16\times 16=256$ candidate phase choices on the two generators finds
exactly $64$ lifts satisfying $x^2=y^3=(xy)^5=1$ and generating a subgroup
of order $60$ meeting $\hat V$ trivially, again always with order-$4$ phases
(gate X5). That no complement can avoid $\mu_4$ in either case is confirmed
by comparing group orders directly (gate X6): the $\mu_2$-phase part of
$\mathrm{Aut}_{\mathcal{F}}^\kappa(A)$ has order $q^2\cdot|O(Q_{\beta_0})\cap SL_2(\kappa)|=4\cdot2=8$
at $q=2$ and $16\cdot 6=96$ at $q=4$, and $|SL_2(\kappa)|=6$, respectively
$60$, does not divide either. This settles the existence question
affirmatively at these two finite values and settles that any splitting
there must use order-$4$ phases; it establishes nothing for $q=2^m$,
$m\geq 3$.
\end{quote}
